% ═══════════════════════════════════════════════════════════════════════════
%  پیوست هـ — تحلیل انتقادی: فروپاشی افغانستان
%  چگونه «نجات‌طلبی» به «رهاسازی» تبدیل شد؟
% ═══════════════════════════════════════════════════════════════════════════
\chapter{تحلیل انتقادی: فروپاشی افغانستان — از دموکراسی‌سازی تا رهاسازی}
\label{app:afghanistan}

\begin{tcolorbox}[mybox, title={چکیده‌ی پیوست}]
این پیوست تحلیلی انتقادی و عمیق از فرایند فروپاشی دولت جمهوری اسلامی افغانستان (۲۰۰۱–۲۰۲۱) ارائه می‌دهد. پرسش محوری آن است: آیا ایالات متحده — به‌جای حمایت از دولت نسبتاً مدرن و دموکراتیک افغانستان — عملاً با مذاکرات غیررسمی و سپس رسمی با طالبان، زمینه‌ی بازگشت آنها را فراهم کرد؟ آیا دستاوردهای بیست‌ساله (حقوق زنان، پارلمان متنوع، رسانه‌های آزاد، قانون اساسی مدرن) قربانی محاسبات ژئوپولیتیک واشنگتن شد؟ تحلیل بر مبنای اسناد رسمی (گزارش‌های \lr{SIGAR})، «اسناد افغانستان» واشنگتن‌پست، پژوهش‌های آکادمیک مستقل و مصاحبه‌ها صورت می‌گیرد.
\end{tcolorbox}

% ─────────────────────────────────────────────────────────────────────────
\section{مقدمه: پارادوکس افغانستان}
\label{sec:app-afg-intro}
% ─────────────────────────────────────────────────────────────────────────

در ۱۵ آگوست ۲۰۲۱، طالبان بدون مقاومت جدی وارد کابل شدند. دولتی که ۲۰ سال، بیش از ۲ تریلیون دلار، و صدها هزار جان برای ساختنش هزینه شده بود، در چند هفته فروپاشید. تصاویر افغان‌هایی که از بدنه‌ی هواپیمای نظامی امریکا آویزان بودند و سقوط می‌کردند، به نماد یکی از بزرگ‌ترین شکست‌های سیاست خارجی معاصر تبدیل شد.%
\footnote{\lr{Whitlock, Craig. \textit{The Afghanistan Papers: A Secret History of the War}. Simon and Schuster, 2021, pp.~xi--xvi.}}

اما آنچه اغلب در روایت رسانه‌ای نادیده گرفته می‌شود، \textbf{ماهیت آنچه از دست رفت} است:

\begin{tcolorbox}[defbox, title={دستاوردهای بیست‌ساله‌ای که از دست رفت}]
\begin{enumerate}[nosep, label=\textbf{\arabic*.}]
  \item \textbf{قانون اساسی ۲۰۰۴:} مدرن‌ترین قانون اساسی در تاریخ افغانستان — شامل تضمین حقوق زنان، آزادی بیان، تفکیک قوا و انتخابات.%
  \footnote{\lr{Constitution of the Islamic Republic of Afghanistan, 2004. Available at \url{https://www.constituteproject.org/}.}}
  \item \textbf{پارلمان با حضور زنان:} در مجلس ۲۰۱۸، ۶۹ زن از ۲۴۹ نماینده (حدود ۲۸٪) — بالاتر از بسیاری کشورهای غربی.%
  \footnote{\lr{Inter-Parliamentary Union. "Women in National Parliaments." 2019.}}
  \item \textbf{رسانه‌های آزاد:} بیش از ۱۷۰ رسانه‌ی فعال (تلویزیون، رادیو، مطبوعات) — شامل رسانه‌های زنان.%
  \footnote{\lr{Reporters Without Borders. "Afghanistan: Media Landscape." 2020.}}
  \item \textbf{آموزش:} حدود ۹ میلیون دانش‌آموز (از جمله ۳.۵ میلیون دختر) در مدارس — در مقایسه با تقریباً صفر دختر در دوره‌ی طالبان.%
  \footnote{\lr{UNICEF Afghanistan. "Education Fact Sheet." 2020.}}
  \item \textbf{جامعه‌ی مدنی:} هزاران سازمان مردم‌نهاد فعال.
  \item \textbf{صنعت فناوری و ارتباطات:} از صفر به ۲۲ میلیون کاربر موبایل.
\end{enumerate}
\end{tcolorbox}

پرسش بنیادین: \textbf{چگونه ممکن است قدرتی که ۲ تریلیون دلار برای ساختن این نظام هزینه کرده، همان قدرت، عامل اصلی فروپاشی آن باشد؟}

% ─────────────────────────────────────────────────────────────────────────
\section{بازخوانی انتقادی روایت رسمی}
\label{sec:app-afg-narrative}
% ─────────────────────────────────────────────────────────────────────────

\subsection{روایت رسمی واشنگتن}
\label{subsec:app-afg-official}

روایت رسمی امریکا (از بوش تا بایدن) مبتنی بر سه ادعاست:

\begin{enumerate}[nosep, label=\textbf{\arabic*.}]
  \item «ما ۲۰ سال به افغان‌ها فرصت دادیم. آنها نتوانستند/نخواستند از آن استفاده کنند.»
  \item «ارتش افغانستان ۳۰۰٬۰۰۰ نیرو داشت اما نخواست بجنگد.»
  \item «مأموریت ما تمام شده بود — القاعده نابود شد.»%
  \footnote{\lr{Biden, Joseph R. "Remarks on Afghanistan." White House, 16 August 2021.}}
\end{enumerate}

\subsection{نقد روایت رسمی: «اسناد افغانستان»}
\label{subsec:app-afg-papers}

\thinker{کریگ ویتلاک} (\lr{Craig Whitlock}) از واشنگتن‌پست، بر اساس هزاران صفحه اسناد محرمانه‌ی افشاشده‌ی \lr{SIGAR} (بازرس ویژه‌ی بازسازی افغانستان)، نشان داد که مقامات امریکایی \textbf{بیست سال آگاهانه به افکار عمومی دروغ گفتند:}

\begin{histQuote}{از «اسناد افغانستان» (واشنگتن‌پست، ۲۰۱۹)}
«اگر مردم امریکا می‌دانستند که واقعاً چه خبر است، جنگ خیلی زودتر تمام شده بود ... ما ارقام [تعداد سربازان افغان، مناطق تحت کنترل دولت، و غیره] را دستکاری می‌کردیم تا نشان دهیم پیشرفت وجود دارد — در حالی که وجود نداشت.»
— \lr{Douglas Lute}, ژنرال سه‌ستاره و هماهنگ‌کننده‌ی جنگ افغانستان در کاخ سفید%
\footnote{\lr{Whitlock, Craig. "At War with the Truth." \textit{The Washington Post}, 9 December 2019. Based on SIGAR, \textit{Lessons Learned} interviews, 2014--2018.}}
\end{histQuote}

\begin{tcolorbox}[enemybox, title={یافته‌های کلیدی «اسناد افغانستان»}]
\begin{enumerate}[nosep, label=\textbf{\arabic*.}]
  \item \textbf{آمار ارتش جعلی بود:} ارتش ۳۰۰٬۰۰۰ نفره‌ی افغانستان بر روی کاغذ وجود داشت. در واقعیت، «سربازان شبح» (\lr{ghost soldiers}) — نام‌هایی که فقط برای دریافت حقوق ثبت شده بودند — بخش بزرگی از آن را تشکیل می‌دادند.%
  \footnote{\lr{SIGAR. "Quarterly Report to Congress." January 2021, pp.~72--78.}}

  \item \textbf{فساد سیستمی:} امریکا آگاهانه فساد دولت افغانستان را نادیده گرفت زیرا «ثبات کوتاه‌مدت» را بر «اصلاحات بلندمدت» ترجیح داد. بخش بزرگی از ۲ تریلیون دلار به جیب پیمانکاران امریکایی، جنگ‌سالاران افغان و مقامات فاسد رفت.%
  \footnote{\lr{Whitlock, 2021, pp.~143--180.}}

  \item \textbf{اهداف متناقض:} امریکا هم‌زمان از دولت افغانستان حمایت و با طالبان مذاکره می‌کرد — بدون آنکه این تناقض را حل کند.

  \item \textbf{پاکستان: متحد یا دشمن؟} امریکا سالانه میلیاردها دلار کمک به پاکستان می‌داد — همان پاکستانی که پناهگاه و حمایت لجستیکی طالبان بود. هیچ تلاش جدی برای توقف حمایت پاکستان از طالبان صورت نگرفت.%
  \footnote{\lr{Coll, Steve. \textit{Directorate S: The C.I.A. and America's Secret Wars in Afghanistan and Pakistan}. Penguin, 2018.}}
\end{enumerate}
\end{tcolorbox}

% ─────────────────────────────────────────────────────────────────────────
\section{توافق دوحه: نقطه‌ی عطف فروپاشی}
\label{sec:app-afg-doha}
% ─────────────────────────────────────────────────────────────────────────

\subsection{ماهیت توافق}

\histref{توافق دوحه}{۲۹ فوریه ۲۰۲۰} میان ایالات متحده و طالبان — \textbf{بدون حضور دولت افغانستان} — نقطه‌ی عطف فروپاشی بود. مفاد اصلی:

\begin{legalNote}{مفاد کلیدی توافق دوحه (فوریه ۲۰۲۰)}
\begin{enumerate}[nosep, label=\textbf{\arabic*.}]
  \item امریکا ظرف ۱۴ ماه تمام نیروهایش را از افغانستان خارج می‌کند.
  \item طالبان تعهد می‌دهد اجازه ندهد قلمروش علیه امریکا استفاده شود.
  \item طالبان و دولت افغانستان وارد «مذاکرات بین‌الافغانی» می‌شوند.
  \item امریکا ۵٬۰۰۰ زندانی طالبان را آزاد می‌کند.
\end{enumerate}
\vspace{4pt}
\textbf{نکته‌ی حیاتی:} هیچ ضمانتی برای حقوق زنان، آزادی مطبوعات، قانون اساسی، یا حفظ ساختار دموکراتیک در توافق وجود نداشت.%
\footnote{\lr{Agreement for Bringing Peace to Afghanistan between the Islamic Emirate of Afghanistan which is not recognized by the United States as a state and is known as the Taliban and the United States of America. Doha, 29 February 2020. Full text available at US Department of State.}}
\end{legalNote}

\subsection{تحلیل انتقادی توافق دوحه}

\begin{tcolorbox}[casebox, title={چرا توافق دوحه «خیانت» به دولت افغانستان بود؟}]

\thinker{کارتر مالکاسیان} (\lr{Carter Malkasian}) — مشاور ارشد ژنرال مک‌کنزی و یکی از بهترین شناسندگان افغانستان — در اثر مهم خود «جنگ امریکا در افغانستان» تحلیل بی‌رحمانه‌ای ارائه می‌دهد:

\begin{enumerate}[nosep, label=\textbf{\arabic*.}]
  \item \textbf{حذف دولت از میز مذاکره:} امریکا مستقیماً با طالبان مذاکره کرد و دولت غنی را کنار گذاشت. این پیام روشنی به جهان فرستاد: امریکا دولت افغانستان را «بازیگر اصلی» نمی‌داند.%
  \footnote{\lr{Malkasian, Carter. \textit{The American War in Afghanistan: A History}. Oxford UP, 2021, pp.~420--458.}}

  \item \textbf{آزادی ۵٬۰۰۰ جنگجوی طالبان:} امریکا دولت غنی را مجبور کرد ۵٬۰۰۰ زندانی طالبان — شامل فرماندهان ارشد — را آزاد کند. بسیاری از آنها مستقیماً به جبهه‌ی جنگ بازگشتند.%
  \footnote{\lr{SIGAR. "Quarterly Report." July 2020, pp.~55--60.}}

  \item \textbf{جدول زمانی بدون شرط:} خروج امریکا \textbf{بدون هیچ شرطی} (مثلاً آتش‌بس، توافق صلح، احترام به قانون اساسی) تعیین شد. طالبان فقط باید تعهد می‌داد «علیه امریکا» اقدام نکند — نه «علیه مردم افغانستان.»

  \item \textbf{از بین‌رفتن روحیه‌ی نظامی:} پس از توافق دوحه، سربازان و افسران افغان فهمیدند که «تنها خواهند ماند.» روحیه‌ی ارتش فروپاشید — پیش از آنکه جنگ واقعی آغاز شود.%
  \footnote{\lr{Malkasian, 2021, pp.~460--475.}}

  \item \textbf{نقش ترامپ و خلیلزاد:} \thinker{زلمی خلیلزاد}، فرستاده‌ی ویژه‌ی ترامپ، «معمار» توافق بود. او خود افغان‌تبار بود اما در عمل \textbf{منافع دولت ترامپ (خروج قبل از انتخابات) را بر منافع مردم افغانستان ترجیح داد.}%
  \footnote{\lr{Filkins, Dexter. "The Surrender." \textit{The New Yorker}, 10 May 2021.}}
\end{enumerate}
\end{tcolorbox}

\subsection{آیا بایدن ادامه‌دهنده بود یا تصمیم‌گیرنده؟}

\begin{table}[htbp]
\centering
\caption{مقایسه‌ی سیاست چهار رئیس‌جمهور امریکا در افغانستان}\label{tab:app-afg-presidents}
\begin{adjustbox}{max width=\linewidth}
\begin{tabular}{
  >{\raggedleft\arraybackslash\bfseries}p{2.5cm}
  >{\raggedleft\arraybackslash}p{3.5cm}
  >{\raggedleft\arraybackslash}p{3.5cm}
  >{\raggedleft\arraybackslash}p{3.5cm}}
\toprule
\rowcolor{PrimaryDeep}
\textcolor{white}{رئیس‌جمهور} &
\textcolor{white}{سیاست اعلامی} &
\textcolor{white}{سیاست واقعی} &
\textcolor{white}{نتیجه} \\
\midrule

بوش (۲۰۰۱–۲۰۰۹) &
نابودی القاعده + دموکراسی‌سازی &
انحراف توجه به عراق (۲۰۰۳)؛ بی‌توجهی به افغانستان &
فرصت اولیه از دست رفت%
\footnote{\lr{Rashid, Ahmed. \textit{Descent into Chaos}. Viking, 2008.}} \\
\rowcolor{NeutralLight}
اوباما (۲۰۰۹–۲۰۱۷) &
«افزایش و خروج» (\lr{Surge and Withdraw}) &
افزایش نیرو (۱۰۰K) سپس کاهش شدید؛ مذاکرات اولیه با طالبان (از ۲۰۱۱) &
نتایج مختلط؛ آغاز روند سازش%
\footnote{\lr{Woodward, Bob. \textit{Obama's Wars}. Simon and Schuster, 2010.}} \\
ترامپ (۲۰۱۷–۲۰۲۱) &
«خاتمه‌ی جنگ‌های بیهوده» &
توافق دوحه؛ کاهش نیرو به ۲٬۵۰۰؛ آزادی زندانیان طالبان &
\textbf{نقطه‌ی عطف فروپاشی}%
\footnote{\lr{Malkasian, 2021, pp.~420--458.}} \\
\rowcolor{NeutralLight}
بایدن (۲۰۲۱) &
اجرای توافق ترامپ + «پایان جنگ بی‌پایان» &
خروج شتاب‌زده بدون برنامه‌ی انتقالی &
فروپاشی فوری \\

\bottomrule
\end{tabular}
\end{adjustbox}
\end{table}

\thinker{بارنت روبین} (\lr{Barnett Rubin})، یکی از برجسته‌ترین کارشناسان افغانستان، استدلال می‌کند که «مشکل بایدن نبود بلکه بیست سال تصمیمات اشتباه بود — اما بایدن آخرین فرصت را نیز از بین برد.»%
\footnote{\lr{Rubin, Barnett R. "Why Did the Afghan Government Collapse?" \textit{Current History}, October 2021.}}

% ─────────────────────────────────────────────────────────────────────────
\section{چرا ارتش ۳۰۰٬۰۰۰ نفره فروپاشید؟}
\label{sec:app-afg-army}
% ─────────────────────────────────────────────────────────────────────────

روایت بایدن — «ارتش ۳۰۰٬۰۰۰ نفره نخواست بجنگد» — ساده‌انگارانه و نادرست است. تحلیل دقیق‌تر:

\begin{tcolorbox}[casebox, title={عوامل واقعی فروپاشی ارتش افغانستان}]
\begin{enumerate}[nosep, label=\textbf{\arabic*.}]
  \item \textbf{«سربازان شبح»:} ارتش ۳۰۰٬۰۰۰ نفره روی کاغذ بود. تخمین‌ها نشان می‌دهد تعداد واقعی نیروهای رزمی حدود ۵۰–۸۰ هزار نفر بود.%
  \footnote{\lr{SIGAR. "Collapse of the ANDSF." Interim Report, February 2022, pp.~24--38.}}

  \item \textbf{وابستگی ساختاری به نیروی هوایی امریکا:} ارتش افغانستان بر اساس الگوی امریکایی ساخته شده بود — الگویی که مبتنی بر پشتیبانی هوایی، اطلاعاتی و لجستیکی بود. وقتی امریکا خارج شد، این ستون فقرات از بین رفت. \textbf{یک ارتش بدون نیروی هوایی عملیاتی، هلیکوپتر، و اطلاعات ماهواره‌ای نمی‌تواند بجنگد — حتی اگر بخواهد.}%
  \footnote{\lr{Malkasian, 2021, pp.~478--490.}}

  \item \textbf{قطع بودجه و لجستیک:} پیمانکاران امریکایی که نگهداری هلیکوپترها و جنگنده‌های افغان را بر عهده داشتند، پس از توافق دوحه خارج شدند. بدون تعمیر و نگهداری، ناوگان هوایی افغانستان عملاً زمین‌گیر شد.%
  \footnote{\lr{SIGAR, 2022, pp.~40--55.}}

  \item \textbf{فروپاشی روحیه:} سربازان افغان ماه‌ها حقوق نگرفته بودند، مهمات کم بود، و پیام سیاسی واشنگتن روشن بود: «شما تنها خواهید ماند.» فرماندهان محلی با طالبان «معامله» کردند — نه از سر خیانت بلکه از سر واقع‌بینی بقایی.%
  \footnote{\lr{Giustozzi, Antonio. \textit{The Taliban at War, 2001--2018}. Oxford UP, 2019, pp.~255--280.}}

  \item \textbf{فساد فرماندهی:} بسیاری از ژنرال‌های افغان فاسد بودند — منصوب‌شده بر اساس وفاداری قومی نه شایستگی. این فساد محصول سیاست امریکا بود که از جنگ‌سالاران به‌عنوان «متحد» استفاده کرد.%
  \footnote{\lr{Whitlock, 2021, pp.~200--235.}}
\end{enumerate}
\end{tcolorbox}

% ─────────────────────────────────────────────────────────────────────────
\section{نقش پاکستان: «متحد» امریکا و «حامی» طالبان}
\label{sec:app-afg-pakistan}
% ─────────────────────────────────────────────────────────────────────────

\begin{tcolorbox}[enemybox, title={پارادوکس پاکستان}]
یکی از بزرگ‌ترین تناقضات جنگ افغانستان، رابطه‌ی سه‌گانه‌ی امریکا-پاکستان-طالبان بود:

\begin{itemize}[nosep, rightmargin=1em]
  \item \textbf{امریکا} سالانه ۱–۲ میلیارد دلار کمک نظامی-اقتصادی به پاکستان می‌داد.%
  \footnote{\lr{Rubin, Barnett R. \textit{Afghanistan from the Cold War through the War on Terror}. Oxford UP, 2013, pp.~325--360.}}
  \item \textbf{پاکستان} (\lr{ISI}) هم‌زمان پناهگاه، آموزش، تسلیحات و اطلاعات به طالبان ارائه می‌داد.%
  \footnote{\lr{Coll, 2018, pp.~301--350.}}
  \item \textbf{امریکا} از این واقعیت آگاه بود اما اقدام جدی نکرد — زیرا:
  \begin{itemize}[nosep]
    \item پاکستان مسیر لجستیک اصلی کمک‌رسانی به افغانستان بود.
    \item پاکستان قدرت هسته‌ای بود و فشار بیش از حد ریسک بی‌ثباتی داشت.
    \item اطلاعات پاکستان برای عملیات ضدتروریسم مورد نیاز بود.
  \end{itemize}
\end{itemize}

\textbf{نتیجه:} امریکا عملاً هم‌زمان یک طرف جنگ (دولت افغانستان) را تأمین مالی و طرف دیگر (طالبان از طریق پاکستان) را غیرمستقیم تقویت می‌کرد. \thinker{استیو کال} (\lr{Steve Coll}) این وضعیت را «شیزوفرنی استراتژیک» نامیده است.%
\footnote{\lr{Coll, 2018, p.~4.}}
\end{tcolorbox}

% ─────────────────────────────────────────────────────────────────────────
\section{آیا «دموکراسی‌سازی» واقعی بود یا ظاهری؟}
\label{sec:app-afg-democracy}
% ─────────────────────────────────────────────────────────────────────────

\subsection{دستاوردهای واقعی}

انکار دستاوردها غیرمنصفانه است. \textbf{برخی تغییرات واقعی بودند:}

\begin{itemize}[nosep, rightmargin=1em]
  \item میلیون‌ها دختر به مدرسه رفتند — این واقعی بود.
  \item رسانه‌های آزاد فعال شدند — این واقعی بود.
  \item زنان وارد پارلمان، دانشگاه، و بازار کار شدند — این واقعی بود.
  \item نسل جدیدی از افغان‌های شهری تربیت شد که ارزش‌های دموکراتیک را درونی کرده بودند — این واقعی بود.
\end{itemize}

\subsection{اما: ساختار شکننده}

\begin{tcolorbox}[defbox, title={چرا دستاوردها شکننده بودند؟}]
\begin{enumerate}[nosep, label=\textbf{\arabic*.}]
  \item \textbf{جزیره‌های مدرنیته:} دستاوردها عمدتاً محدود به شهرهای بزرگ (کابل، هرات، مزارشریف) بودند. در ۷۰٪ افغانستان (مناطق روستایی)، تغییر چندانی رخ نداده بود. \thinker{توماس بارفیلد} (\lr{Thomas Barfield}) تأکید می‌کند که «افغانستان شهری و افغانستان روستایی دو جهان متفاوت بودند.»%
  \footnote{\lr{Barfield, Thomas. \textit{Afghanistan: A Cultural and Political History}. Princeton UP, 2010, pp.~318--340.}}

  \item \textbf{دموکراسی وارداتی:} نظام سیاسی بر اساس الگوی غربی طراحی شده بود — ریاست‌جمهوری متمرکز (نه فدرالی)، انتخابات حزبی (در جامعه‌ی قبیله‌ای) — که با واقعیت‌های اجتماعی افغانستان سازگار نبود.%
  \footnote{\lr{Suhrke, Astri. \textit{When More Is Less: The International Project in Afghanistan}. Columbia UP, 2011.}}

  \item \textbf{وابستگی مالی مطلق:} دولت افغانستان حدود ۷۵٪ بودجه‌اش را از کمک‌های خارجی دریافت می‌کرد. «دولتی که پول مردمش را خرج نمی‌کند، به مردمش پاسخ‌گو نیست.»%
  \footnote{\lr{Goodhand, Jonathan. "Corrupting or Consolidating the Peace? The Drugs Economy and Post-Conflict Peacebuilding in Afghanistan." \textit{International Peacekeeping} 15, no.~3 (2008): 405--423.}}

  \item \textbf{فمینیسم ابزاری:} \thinker{آنیکا رابو} (\lr{Annika Rabo}) و دیگر پژوهشگران فمینیست استدلال کرده‌اند که «حقوق زنان» در افغانستان \textbf{هم واقعی و هم ابزاری} بود: واقعی زیرا زنان افغان واقعاً آزادتر شدند؛ ابزاری زیرا واشنگتن از آن برای مشروعیت‌بخشی به جنگ استفاده کرد.%
  \footnote{\lr{Abu-Lughod, Lila. "Do Muslim Women Need Saving?" \textit{American Anthropologist} 104, no.~3 (2002): 783--790.}}
\end{enumerate}
\end{tcolorbox}

% ─────────────────────────────────────────────────────────────────────────
\section{نمودار علّی: فرایند فروپاشی}
\label{sec:app-afg-causal}
% ─────────────────────────────────────────────────────────────────────────

\begin{figure}[htbp]
\centering
\begin{tikzpicture}[
  every node/.style={font=\small},
  cause/.style={
    draw=AccentRed, fill=BgFail!50, rounded corners=4pt,
    text width=3.5cm, align=center, minimum height=1.2cm
  },
  process/.style={
    draw=AccentAmber, fill=BgWarm, rounded corners=4pt,
    text width=3.5cm, align=center, minimum height=1.2cm
  },
  result/.style={
    draw=PrimaryDeep, fill=BgCool, rounded corners=4pt,
    text width=4cm, align=center, minimum height=1.5cm,
    font=\small\bfseries
  },
  arrow/.style={->, >=Stealth, thick, NeutralDark},
]

% ردیف بالا: علل ساختاری
\node[cause] (c1) at (0,8) {\rl{فساد سیستمی دولت}};
\node[cause] (c2) at (5,8) {\rl{وابستگی مالی مطلق}};
\node[cause] (c3) at (10,8) {\rl{ارتش «کاغذی»}};

% ردیف دوم: علل خارجی
\node[cause] (c4) at (0,5.5) {\rl{حمایت پاکستان از طالبان}};
\node[cause] (c5) at (5,5.5) {\rl{انحراف به عراق (۲۰۰۳)}};
\node[cause] (c6) at (10,5.5) {\rl{خستگی امریکا از جنگ}};

% نقطه‌ی عطف
\node[process, fill=AccentAmber!30, text width=5cm] (doha) at (5,3)
  {\rl{\textbf{توافق دوحه (۲۰۲۰)}\\%
   حذف دولت از مذاکره\\%
   آزادی ۵۰۰۰ طالبان\\%
   بدون شرط حقوق بشر}};

% فرایند فروپاشی
\node[process] (p1) at (0,0.5) {\rl{فروپاشی روحیه‌ی ارتش}};
\node[process] (p2) at (5,0.5) {\rl{معاملات محلی با طالبان}};
\node[process] (p3) at (10,0.5) {\rl{خروج شتاب‌زده (۲۰۲۱)}};

% نتیجه
\node[result, fill=AccentRed!15] (res) at (5,-2)
  {\rl{\Large سقوط کابل}\\%
   \rl{۱۵ آگوست ۲۰۲۱}\\%
   \rl{بازگشت طالبان}};

% فلش‌ها
\draw[arrow] (c1) -- (doha);
\draw[arrow] (c2) -- (doha);
\draw[arrow] (c3) -- (doha);
\draw[arrow] (c4) -- (doha);
\draw[arrow] (c5) -- (doha);
\draw[arrow] (c6) -- (doha);
\draw[arrow] (doha) -- (p1);
\draw[arrow] (doha) -- (p2);
\draw[arrow] (doha) -- (p3);
\draw[arrow] (p1) -- (res);
\draw[arrow] (p2) -- (res);
\draw[arrow] (p3) -- (res);

\end{tikzpicture}
\caption{نمودار علّی فروپاشی دولت افغانستان}
\label{fig:app-afg-causal}
\end{figure}

% ─────────────────────────────────────────────────────────────────────────
\section{نجات‌طلبی در افغانستان پس از سقوط}
\label{sec:app-afg-post-fall}
% ─────────────────────────────────────────────────────────────────────────

\subsection{آیا گروهی در افغانستان «نجات‌طلب» است؟}

پس از سقوط کابل، چندین گروه خواستار مداخله‌ی خارجی (یا بازگشت مداخله) شده‌اند:

\begin{longtable}{
  >{\raggedleft\arraybackslash\bfseries}p{2.5cm}
  >{\raggedleft\arraybackslash}p{3cm}
  >{\raggedleft\arraybackslash}p{3.5cm}
  >{\raggedleft\arraybackslash}p{4cm}}
\caption{گروه‌های نجات‌طلب در افغانستان پس از ۲۰۲۱}\label{tab:app-afg-seekers}\\
\toprule
\rowcolor{PrimaryDeep}
\textcolor{white}{\textbf{گروه}} &
\textcolor{white}{\textbf{خواسته}} &
\textcolor{white}{\textbf{از چه کسی؟}} &
\textcolor{white}{\textbf{واقع‌بینی}} \\
\midrule
\endfirsthead
\bottomrule
\endlastfoot

زنان فعال مدنی &
فشار بین‌المللی بر طالبان &
سازمان ملل، غرب &
بسیار محدود — غرب دیگر اراده‌ی مداخله ندارد%
\footnote{\lr{Human Rights Watch. "Afghanistan: Taliban Deprive Women of Livelihoods, Identity." January 2022.}} \\
\rowcolor{NeutralLight}
جبهه‌ی مقاومت ملی (پنجشیر) &
حمایت نظامی/تسلیحاتی &
تاجیکستان، هند، غرب &
بسیار ضعیف — فقدان حمایت جدی خارجی%
\footnote{\lr{Giustozzi, Antonio. "The Resistance After the Taliban Takeover." \textit{RUSI Journal} 167, no.~1 (2022): 24--33.}} \\
دیاسپورای افغان &
عدم شناسایی طالبان + تحریم &
اتحادیه‌ی اروپا، امریکا &
متوسط — اما بدون اهرم فشار مؤثر \\
\rowcolor{NeutralLight}
هزاره‌ها &
حمایت در برابر سرکوب &
ایران (محدود)، سازمان ملل &
بسیار ضعیف — هزاره‌ها آسیب‌پذیرترین گروه هستند%
\footnote{\lr{Amnesty International. "Afghanistan: Hazara Targeted in Escalating Attacks." October 2021.}} \\
نخبگان پیشین (تکنوکرات‌ها) &
حکومت فراگیر + بازگشت &
امریکا، قطر &
ناچیز — طالبان حاضر به مشارکت واقعی نیست \\

\end{longtable}

\subsection{تلخ‌ترین درس: نجات‌طلبی بدون ناجی}

\begin{tcolorbox}[enemybox, title={وضعیت کنونی: نجات‌طلبی بدون ناجی}]
میلیون‌ها افغان — به‌ویژه زنان، اقلیت‌ها و نسل جوان شهری — \textbf{خواهان نجات هستند اما ناجی‌ای وجود ندارد:}

\begin{itemize}[nosep, rightmargin=1em]
  \item \textbf{امریکا:} دیگر اراده‌ی بازگشت ندارد. اوکراین اولویت است.
  \item \textbf{اروپا:} نگران مهاجرت است نه حقوق بشر افغان‌ها.
  \item \textbf{سازمان ملل:} ناتوان و بی‌اثر.
  \item \textbf{ایران:} خود از حقوق بشر فاصله دارد و با طالبان رابطه‌ی عملی دارد.
  \item \textbf{پاکستان:} حامی طالبان است.
  \item \textbf{چین و روسیه:} خواهان ثبات‌اند نه دموکراسی.
\end{itemize}

\textbf{نتیجه:} مردم افغانستان در بدترین حالت ممکن قرار دارند: \concept{وابسته‌سازی‌شده و سپس رهاشده}. بیست سال مداخله، جامعه‌ی سنتی را تا حدی مدرن کرد اما سپس بدون هیچ شبکه‌ی ایمنی، آن را به نظام قرون‌وسطایی بازگرداند.
\end{tcolorbox}

% ─────────────────────────────────────────────────────────────────────────
\section{افغانستان از منظر الگوی «مست گیرند...»}
\label{sec:app-afg-paradigm}
% ─────────────────────────────────────────────────────────────────────────

\begin{tcolorbox}[wavebox, title={تحلیل نهایی: افغانستان و قاعده‌ی «مست گیرند، در شهر هر آنکه هست گیرند»}]

افغانستان \textbf{کامل‌ترین و تلخ‌ترین نمونه‌ی تاریخی} قاعده‌ی «مست گیرند...» است — اما با یک چرخش: مداخله‌گر نه‌فقط «مست نگرفت» بلکه \textbf{پس از ۲۰ سال، همه را رها کرد:}

\begin{enumerate}[nosep, label=\textbf{\arabic*.}]
  \item \textbf{فاز اول (۲۰۰۱–۲۰۰۳): «آزادسازی» واقعی.} سقوط طالبان، آزادی نسبی، امید. این فاز نسبتاً موفق بود.

  \item \textbf{فاز دوم (۲۰۰۳–۲۰۰۹): انحراف و بی‌توجهی.} توجه به عراق. افغانستان فراموش شد. طالبان بازسازی شدند.

  \item \textbf{فاز سوم (۲۰۰۹–۲۰۱۴): «افزایش» ناکام.} اوباما ۱۰۰٬۰۰۰ سرباز فرستاد اما استراتژی مشخصی نداشت. فساد عمیق‌تر شد.

  \item \textbf{فاز چهارم (۲۰۱۴–۲۰۲۰): «عقب‌نشینی تدریجی».} کاهش نیرو، مذاکره با طالبان، تضعیف دولت.

  \item \textbf{فاز پنجم (۲۰۲۰–۲۰۲۱): «رهاسازی».} توافق دوحه + خروج شتاب‌زده = فروپاشی.
\end{enumerate}

\textbf{درس بنیادین:}

قاعده‌ی «مست گیرند...» را باید بازنویسی کرد:

\begin{center}
\textbf{\large «مداخله‌ی خارجی نه‌فقط ذاتاً گسترش‌یابنده است، بلکه ذاتاً بازگشت‌پذیر نیز هست.\\[4pt]
آنچه مداخله‌گر می‌سازد، مداخله‌گر می‌تواند ویران کند — \\[4pt]
نه با حمله بلکه صرفاً با رفتن.»}
\end{center}

این «ویرانی با رفتن» (\lr{destruction by departure}) مفهوم جدیدی است که افغانستان به ادبیات مداخلات خارجی اضافه کرده: \concept{خطرناک‌ترین مداخله آنی نیست که با خشونت پایان یابد — بلکه آنی است که پس از بیست سال وابسته‌سازی، ناگهان متوقف شود.}%
\footnote{\lr{Malkasian, 2021, p.~498: "The greatest failure was not the intervention itself, but the manner of its ending."}}
\end{enumerate}
\end{tcolorbox}

% ─────────────────────────────────────────────────────────────────────────
\section*{نتیجه‌گیری: افغانستان به‌عنوان آزمون نهایی نجات‌طلبی}%

\label{sec:app-afg-conclusion}
% ─────────────────────────────────────────────────────────────────────────

\begin{tcolorbox}[wavebox, title={یافته‌های نهایی}]
\begin{enumerate}[nosep, label=\textbf{\arabic*.}]
  \item \textbf{درک شما (خواننده) عمدتاً درست است — اما نیاز به ظرافت دارد:} امریکا واقعاً با مذاکره‌ی مستقیم با طالبان (و حذف دولت)، بازگشت آنها را تسهیل کرد. اما علت‌ها پیچیده‌تر از «توطئه» هستند: خستگی جنگ، شکست دولت‌سازی، فساد، نقش پاکستان و ناتوانی ساختاری.

  \item \textbf{فروپاشی نه «ناگهانی» بلکه «تدریجی» بود:} بیست سال تصمیمات اشتباه (انحراف به عراق، تحمل فساد، نادیده‌گرفتن پاکستان) زیربنای فروپاشی را ساخت. توافق دوحه «آخرین میخ» بود نه «تنها علت.»

  \item \textbf{افغانستان نمونه‌ی \concept{مداخله‌ی «خوب» است که بد تمام شد}:} مداخله‌ی ۲۰۰۱ دلایل مشروع‌تری نسبت به عراق ۲۰۰۳ داشت (دفاع مشروع، مجوز شورای امنیت). دستاوردهای واقعی ایجاد شد. اما وابسته‌سازی ساختاری + بی‌مسئولیت خروج = فاجعه.

  \item \textbf{قربانیان اصلی مردم عادی هستند:} زنان، دخترانی که به مدرسه رفتند و حالا در خانه زندانی‌اند؛ خبرنگارانی که صدا داشتند و حالا ساکت‌اند؛ هزاره‌هایی که از نسل‌کشی در امان بودند و حالا دوباره هدف هستند.

  \item \textbf{هیچ‌کس پاسخ‌گو نیست:} نه بوش (که جنگ را آغاز کرد)، نه اوباما (که آن را بی‌سرانجام رها کرد)، نه ترامپ (که با طالبان معامله کرد)، نه بایدن (که خروج فاجعه‌بار را اجرا کرد). «پاسخ‌گونبودن مداخله‌گر» یکی از بزرگ‌ترین نقص‌های نظام بین‌الملل است.

  \item \textbf{پیام برای هر ملت نجات‌طلبی:}
  \begin{center}
  \textit{«اگر آزادی‌ات را مدیون ارتش بیگانه باشی،\\
  آزادی‌ات به اندازه‌ی حضور آن ارتش دوام خواهد داشت — \\
  نه یک روز بیشتر.»}
  \end{center}
\end{enumerate}
\end{tcolorbox}

\cleardoublepage